% ============================================================
%  capitulo5_metodologia.tex
%  Capítulo 5: Metodología Integral para la Normalización
%              de Redes Eléctricas con AGPE en el Marco
%              del PRONE
% ============================================================

\chapter{Metodología Integral para la Normalización de Redes Eléctricas con AGPE en el Marco del PRONE}
\label{chap:metodologia}

Este capítulo es el producto central de la monografía. Todo lo que se
desarrolló en los capítulos anteriores —el diagnóstico de los BS, la
revisión normativa, los casos de estudio y el análisis técnico y
financiero— converge aquí en una propuesta metodológica concreta para
que los OR puedan formular y presentar proyectos de normalización con
AGPE que sean elegibles ante el MME.

La metodología no es un protocolo rígido. Es una hoja de ruta adaptable
que organiza los componentes técnicos, normativos, financieros, sociales
e institucionales del proceso en una secuencia lógica, con actores
definidos, entregables claros y criterios de avance entre fases.


% ─────────────────────────────────────────────────────────────
\section{Identificación de Actores y sus Roles en el Proceso}
\label{sec:actores}
% ─────────────────────────────────────────────────────────────

Normalizar un BS no es algo que un OR pueda hacer solo. Lo demostró
Popayán, que requirió la coordinación de más de diez entidades para
normalizar 41 asentamientos \parencite{SerranoRueda2019}. La
articulación de actores no es un requisito de forma: es la condición
que determina si el proyecto puede superar los bloqueos que
inevitablemente aparecen en el camino.

Los actores del proceso se agrupan en cinco categorías:

\subsection{Actores públicos nacionales}
\label{subsec:actores_nacionales}

\textbf{Ministerio de Minas y Energía (MME):} Administra el PRONE,
define los lineamientos de cada convocatoria, los requisitos de
elegibilidad y los criterios de priorización. Expide los actos
administrativos de adjudicación y supervisa la ejecución. Es el
destinatario final del expediente del proyecto.

\textbf{CAPRONE:} El comité colegiado que evalúa y aprueba los proyectos
en cada convocatoria. Sus actas son el registro oficial de lo que se
adjudica y a quién.

\textbf{CREG:} Define el marco regulatorio para la distribución y para
la integración de AGPE en el SIN, incluyendo los requisitos de conexión,
medición bidireccional y gestión de excedentes. La Resolución 174 de 2021
es la norma vigente \parencite{ResolucionCREG174_2021}.

\textbf{UPME:} Define la planeación energética nacional y publica los
datos de consumo de referencia por piso térmico. También fija los límites
de potencia para los sistemas AGPE \parencite{UPME2025}.

\textbf{SSPD:} Es la administradora del SUI, donde se centralizan los
reportes realizados por los prestadores de servicios públicos del SIN a
través de formatos previamente estructurados. Esta información no solo
permite la vigilancia de las entidades en la cadena de valor, sino que
representa el insumo para informes, estadísticas, investigaciones y
control social, entre otros. Por tanto, al ser el sistema definido para
ello, el MME se basa en el registro de los formatos S4, TC1 y TC2 para
la asignación del subsidio FOES y la habilitación de los barrios
subnormales en sus convocatorias.

\subsection{Actores territoriales}
\label{subsec:actores_territoriales}

\textbf{Alcaldías municipales y distritales:} El actor territorial más
crítico. Sin la certificación del alcalde, ningún proyecto puede
presentarse al PRONE. Punto. No importa cuánto tiempo lleve el barrio en
subnormalidad ni qué tan bien esté diseñado el proyecto técnico. Es a
través de su Plan de Ordenamiento Territorial (POT), Plan Básico de
Ordenamiento Territorial (PBOT) o Esquema de Ordenamiento Territorial
(EOT) que se reconoce legalmente la subnormalidad a través de la
certificación. Las alcaldías también son actores clave en la gestión
social previa y en la articulación con la comunidad durante las obras.
Adicionalmente, una vez las redes de los BS sean aptas, la entidad
territorial es la encargada de prestar el alumbrado público, garantizando
continuidad, calidad del servicio y los niveles de cobertura adecuados
conforme a lo estipulado en el artículo 4 del Decreto 943 de 2018
\parencite{Decreto943_2018}.

\textbf{Gobernaciones:} Pueden articular entre el nivel nacional y el
municipal, especialmente en departamentos con alta concentración de BS
y débil capacidad institucional municipal.

\subsection{Actores operativos del sector eléctrico}
\label{subsec:actores_operativos}

\textbf{Operadores de Red (OR):} Tienen la mayor responsabilidad del
proceso. Deben identificar y diagnosticar los BS en su área, gestionar
la certificación ante los entes territoriales, formular y presentar
el proyecto, contratar y supervisar las obras, y asumir la AOM de los
activos de la red una vez terminadas las obras. En proyectos con AGPE,
también asumen la responsabilidad técnica de la conexión al SIN, salvo
que los activos de generación se cedan a una Comunidad Energética
\parencite{Decreto1023_2024}.

\textbf{Comercializadores:} En los BS, el OR frecuentemente actúa
también como comercializador. Cuando opera un comercializador distinto,
se introduce una capa adicional de coordinación para el manejo de los
esquemas diferenciales de facturación.

\subsection{Actores privados de soporte}
\label{subsec:actores_privados}

\textbf{Contratistas y constructores:} Ejecutan las obras bajo
supervisión del OR. Para trabajar en BS no basta con experiencia técnica
eléctrica: los trabajadores de obra son el primer contacto visible del
proyecto con la comunidad, y eso importa.

\textbf{Proveedores de equipos:} Suministran materiales de red y equipos
del sistema AGPE. En la región Caribe, la disponibilidad local de
proveedores con garantía y postventa es un criterio de selección
relevante que no siempre se considera en la formulación.

\textbf{Firmas consultoras:} Apoyan al OR en la formulación, el diseño
técnico, el presupuesto y la preparación del expediente. Son
especialmente valiosas en los OR con menor capacidad interna de
estructuración, que son precisamente los que históricamente han tenido
más problemas para presentar proyectos elegibles.

\subsection{Actores sociales y comunitarios}
\label{subsec:actores_comunitarios}

\textbf{Juntas de Acción Comunal (JAC):} El interlocutor comunitario
más directo. Su respaldo al proceso puede ser la diferencia entre una
obra que avanza y una que la comunidad bloquea. Cuando la JAC asume el
rol de suscriptor comunitario durante la transición, esa figura debe
manejarse con cuidado: la JAC no tiene herramientas legales para exigir
el pago individual a los usuarios.

\textbf{Líderes sociales e informales:} En los BS hay liderazgos con
más influencia real que la JAC formal. Identificarlos y vincularlos
desde el inicio del proceso es una de las acciones con mayor retorno
en la gestión social.

\textbf{Comunidades Energéticas:} Figura reglamentada por el Decreto
2236 de 2023 que permite a un grupo de usuarios asumir el carácter
legal de prestador de servicio público y gestionar los activos AGPE
construidos con recursos del PRONE. La Comunidad Energética puede estar
facultada para la gestión de los activos del sistema de AGPE y para
suscribir contratos o convenios asociativos para desarrollar las
actividades de generación, comercialización, y uso eficiente de la
energía. Su constitución requiere un proceso de formación comunitaria
que debe iniciarse con suficiente anticipación a la finalización de las
obras \parencite{Decreto2236_2023}.

La Tabla~\ref{tab:actores_roles} sintetiza los actores, su fase de
intervención y su función principal.

\begin{table}[H]
\centering
\caption{Actores del proceso de normalización con AGPE, fases de intervención y función principal}
\label{tab:actores_roles}
\begin{tabularx}{\textwidth}{>{\RaggedRight}p{2.2cm}
                             >{\RaggedRight}p{3.8cm}
                             >{\RaggedRight}p{2.8cm}
                             >{\RaggedRight\arraybackslash}X}
\toprule
\textbf{Tipo} & \textbf{Entidad} & \textbf{Fase} & \textbf{Función principal} \\
\midrule
\rowcolor{grisclaro}
Públicos nacionales & MME, CAPRONE, CREG, UPME, SSPD & Planeación, regulación y aprobación & Administración del PRONE, definición de políticas y normas, aprobación de proyectos y supervisión. \\
Territoriales & Alcaldías, gobernaciones & Diagnóstico y gestión & Certificación del BS, articulación comunitaria, acompañamiento social y soporte institucional local. \\
\rowcolor{grisclaro}
Operativos & OR, comercializadores & Planeación, ejecución y operación & Diagnóstico técnico, diseño, formulación, construcción, puesta en servicio y AOM. \\
Privados de soporte & Contratistas, proveedores, consultoras & Ejecución & Materiales, mano de obra calificada y consultoría técnica. \\
\rowcolor{grisclaro}
Sociales y comunitarios & JAC, líderes sociales, Comunidades Energéticas & Diagnóstico, ejecución y seguimiento & Socialización, apropiación del proyecto, gestión comunitaria y AOM de activos AGPE. \\
\bottomrule
\multicolumn{4}{l}{\footnotesize\textit{Fuente: Elaboración propia.}}
\end{tabularx}
\end{table}


% ─────────────────────────────────────────────────────────────
\section{Requisitos para la Financiación a través del PRONE}
\label{sec:requisitos}
% ─────────────────────────────────────────────────────────────

En 2024, ningún proyecto cumplió los requisitos del PRONE
\parencite{MME2025}. Ese resultado no es anecdótico: es la señal más
clara de que muchos OR no dimensionan bien la complejidad de lo que se
les pide. Los requisitos abarcan cuatro dimensiones, y el incumplimiento
de cualquiera de ellas es causal de descalificación.

\subsection{Requisitos técnicos}
\label{subsec:req_tecnicos}

\textbf{Para la red de normalización:}
\begin{itemize}
    \item Diseño eléctrico completo conforme al RETIE (Resolución MME
    40117 de 2024), incluyendo análisis de riesgos, cálculos de carga,
    selección de conductores, coordinación de protecciones, sistema de
    puesta a tierra, diagramas unifilares y planos de construcción.
    \item Los gastos elegibles se limitan a redes de distribución,
    transformadores, acometidas y medidores. El desmonte del material
    existente no puede superar el 3\% del valor total del proyecto
    \parencite{DUR1073_2015}.
    \item Cumplimiento del Código Eléctrico Colombiano (NTC 2050) y de
    las normas técnicas del OR.
    \item Diseños de ingeniería, memorias de cálculo y planos (eléctricos
    y civiles), aportados a título gratuito al MME.
    \item Cronograma proyectado para la ejecución del proyecto y las
    actividades requeridas.
    \item Plan donde se determinen las estrategias para garantizar la
    sostenibilidad del proyecto a nivel social, ambiental y económico.
\end{itemize}

\textbf{Para el sistema AGPE:}
\begin{itemize}
    \item Potencia máxima instalada: 1~MW, conforme a la Resolución
    UPME 281 de 2015 \parencite{ResolucionUPME281_2015}.
    \item La potencia del AGPE no puede superar el 50\% de la capacidad
    del transformador o circuito al que se conecta, conforme a la
    CREG 174 de 2021 \parencite{ResolucionCREG174_2021}.
    \item Medición bidireccional obligatoria.
    \item El diseño debe incluir los estudios de recurso solar,
    dimensionamiento del sistema, verificación del límite del 50\% y
    compatibilidad con la red del OR.
    \item Si los activos se ceden a una Comunidad Energética: acto de
    constitución conforme al Decreto 2236 de 2023 y plan de AOM
    aprobado por el OR.
    \item Diagramas unifilares existentes y proyectados, incluyendo
    sistema de infraestructura de medición avanzada (AMI).
    \item Análisis de alternativas del sistema AGPE.
\end{itemize}

\subsection{Requisitos sociales}
\label{subsec:req_sociales}

Los 20 expertos consultados dieron 5 sobre 5 al acompañamiento social.
No hubo una sola excepción. Eso significa que el componente social no
es opcional en la formulación del proyecto: es parte estructural del
expediente. Los requisitos concretos son:

\begin{itemize}
    \item \textbf{Diagnóstico social previo:} Caracterización
    socioeconómica y cultural del BS, incluyendo liderazgos, historias
    previas de intervención, presencia étnica y nivel de resistencia
    esperado.

    \item \textbf{Plan de gestión social:} Documento con cronograma,
    metodología, responsables y presupuesto para las acciones de
    socialización antes, durante y después de las obras.

    \item \textbf{Acuerdo con el suscriptor comunitario o JAC} sobre el
    esquema de facturación diferencial durante la transición.

    \item \textbf{Plan de eficiencia energética:} Acciones para reducir
    el consumo de los hogares post-normalización, de manera que el
    impacto de la factura sea manejable. En proyectos con AGPE, debe
    incluir información sobre los beneficios esperados del sistema solar
    en la factura del usuario.

    \item \textbf{Posición clara sobre las deudas históricas:} La
    experiencia de Popayán es contundente — sin condonación de deudas,
    no hubo normalización sostenible \parencite{SerranoRueda2019}.
\end{itemize}

\subsection{Requisitos económicos}
\label{subsec:req_economicos}

\begin{itemize}
    \item \textbf{Presupuesto detallado:} Desagregado por componentes
    (red MT, transformadores, red BT, medidores, acometidas, AGPE,
    gestión social, interventoría), con cantidades de obra, precios
    unitarios verificables y valores totales. Incluye presupuesto y
    costos del proyecto a nivel global y unitario.

    \item \textbf{Proyecto formulado en la Metodología General Ajustada
    (MGA):} Herramienta diseñada por el DNP para la formulación y
    evaluación de proyectos de inversión pública, requerida para la
    presentación del expediente.

    \item \textbf{Plan de inversiones quinquenal:} Con recursos del OR,
    demostrando la capacidad de inversión y el compromiso financiero
    con el proyecto.

    \item \textbf{Costo por usuario:} El CAPRONE evalúa si la relación
    entre el valor del proyecto y el número de familias beneficiadas es
    coherente con el histórico del programa. La referencia disponible
    es \$1.093.052 por usuario en proyectos sin AGPE
    \parencite{CabralesSanjuan2015}, actualizable por inflación y
    complementada con el costo adicional del componente solar.

    \item \textbf{Análisis de impacto tarifario:} Estimación de la
    factura mensual esperada por usuario normalizado con y sin AGPE,
    comparada con el subsidio FOES. Debe mostrar que la tarifa
    resultante es compatible con la capacidad de pago del BS.

    \item \textbf{Proyección de recaudo:} Estimación conservadora del
    recaudo esperado en el primer, segundo y quinto año
    post-normalización, teniendo en cuenta la curva de adaptación de
    los usuarios a la facturación individual.
\end{itemize}

\subsection{Requisitos legales y administrativos}
\label{subsec:req_legales}

Son los más estrictamente verificados por el CAPRONE y los que con más
frecuencia generan descalificación:

\begin{itemize}
    \item \textbf{Certificación de BS vigente,} expedida por el alcalde
    conforme al artículo 2.2.3.1.2 del DUR 1073 de 2015, actualizada y
    referenciando con precisión el barrio objeto del proyecto
    \parencite{DUR1073_2015}.

    \item \textbf{Carta de presentación} donde se identifiquen los datos
    generales del proyecto incluyendo el recurso solicitado.

    \item \textbf{Garantía de seriedad y/o cumplimiento de la oferta,}
    conforme a lo establecido en los términos de la convocatoria.

    \item \textbf{Acuerdo OR--comercializador} si son diferentes, para
    atender usuarios y comercializar excedentes del sistema AGPE.

    \item \textbf{Certificado de registro de los BS en el SUI,} que
    acredite que los barrios objeto del proyecto se encuentran reportados
    ante la SSPD.

    \item \textbf{Carta del ente territorial comprometiéndose con el
    alumbrado público} una vez se normalicen las redes, conforme al
    Decreto 943 de 2018 \parencite{Decreto943_2018}.

    \item \textbf{Certificación de la entidad territorial} que conste que
    el barrio es BS y defina la estratificación socioeconómica aplicable.

    \item \textbf{Acreditación del OR ante el PRONE,} incluyendo la
    condición de prestador del servicio en el área del proyecto y el
    cumplimiento de sus obligaciones con el SUI.

    \item \textbf{Estudio de títulos del terreno} para verificar que no
    hay restricciones legales que impidan ejecutar las obras.

    \item \textbf{Permisos y licencias ambientales} cuando la normativa
    vigente los exija para la instalación de infraestructura de
    distribución o de generación solar.

    \item \textbf{Constitución de la Comunidad Energética} (si aplica),
    conforme al Decreto 2236 de 2023, como parte del expediente
    \parencite{Decreto2236_2023}.

    \item \textbf{Pólizas de cumplimiento} que garanticen las
    obligaciones de ejecución y entrega ante el MME.
\end{itemize}

\subsection{Requisitos adicionales de convocatorias recientes}
\label{subsec:req_adicionales}

Las convocatorias más recientes del PRONE han incorporado requisitos
adicionales que complementan las dimensiones anteriores y que deben
tenerse en cuenta en la formulación del expediente:

\begin{itemize}
    \item Suscripción del Plan de Acción de Responsabilidad Social del
    MME, que incluye las obligaciones asociadas a la conformación de
    Comunidades Energéticas cuando el proyecto incorpore AGPE.

    \item Restricción de que el costo de la interventoría no supere el
    7\% de los costos directos del proyecto.

    \item Garantizar como mínimo el consumo básico de subsistencia para
    los usuarios beneficiados.

    \item Listado de usuarios beneficiados, con identificación y
    ubicación de cada vivienda.

    \item Certificación del OR sobre contratación de al menos 50\% de
    personal local, apoyo a redes internas, punto de conexión del
    sistema AGPE y realización de socializaciones con la comunidad.

    \item Documento de cofinanciación, si aplica, indicando la fuente y
    el monto de los recursos complementarios.

    \item Declaración de cumplimiento de los profesionales responsables
    del diseño y la ejecución del proyecto.

    \item Análisis de riesgo del proyecto, que identifique las amenazas
    técnicas, sociales, ambientales y financieras y las medidas de
    mitigación correspondientes.
\end{itemize}


% ─────────────────────────────────────────────────────────────
\section{Hoja de Ruta Metodológica: Fases del Proceso}
\label{sec:hoja_de_ruta}
% ─────────────────────────────────────────────────────────────

La metodología se organiza en cuatro fases secuenciales y una fase
transversal de gestión social y comunitaria que acompaña todo el proceso.
Cada fase tiene un objetivo claro, actores responsables definidos,
entregables mínimos y un criterio de avance que determina cuándo puede
pasarse a la siguiente. No se puede saltar una fase. Intentarlo es
exactamente lo que explica por qué tantos proyectos llegan incompletos
al CAPRONE.

\begin{figure}[H]
\centering
%\includegraphics[width=\textwidth]{figura_metodologia_fases.png}
\caption{Fases de la metodología integral para normalización de redes eléctricas con AGPE en el marco del PRONE}
\label{fig:fases_metodologia}
\footnotesize\textit{Fuente: Elaboración propia.}
\end{figure}

\subsection{Fase 1 -- Diagnóstico integral del territorio}
\label{subsec:fase1}

\textbf{Objetivo:} Conocer el BS en todas sus dimensiones antes de
diseñar cualquier solución.

\textbf{Responsables:} OR (liderazgo), alcaldía (certificación y datos
territoriales), trabajo social (diagnóstico social), consultora técnica
si aplica.

\textbf{Actividades:}

\begin{enumerate}[label=\textbf{\arabic*.}, leftmargin=1.8cm, itemsep=6pt]

    \item \textbf{Verificar la certificación del BS.} Si no existe o
    está vencida, iniciar las gestiones ante la alcaldía de inmediato.
    Esta gestión puede tardar meses y debe arrancar con al menos seis
    meses de anticipación respecto a la fecha esperada de la convocatoria.

    \item \textbf{Censo de usuarios y levantamiento técnico de la red:}
    Número real de viviendas, inventario de infraestructura existente
    (postes, conductores, transformadores, macromedidores), estimación
    del consumo actual y caracterización del estado técnico de las redes
    internas.

    \item \textbf{Diagnóstico socioeconómico y cultural:} Condiciones
    de ingresos, estratificación, capacidad de pago, presencia de
    comunidades étnicas, liderazgos formales e informales, y experiencias
    previas del OR en ese territorio. Este último punto es especialmente
    importante: si el OR ya intervino antes y la comunidad quedó
    inconforme, eso define el punto de partida de la gestión social.

    \item \textbf{Recurso solar y condiciones climáticas:} Piso térmico
    del municipio conforme a la Resolución 194 de 2025 del Ministerio
    de Vivienda, datos de irradiación (GHI y HPS) del sitio desde el
    Atlas de Radiación Solar del IDEAM, NASA POWER o Global Solar Atlas,
    y consumo real de referencia a partir de la plataforma O3 de la SSPD
    \parencite{SSPD2025, UPME2025}. No usar los estándares del PIEC
    para dimensionar: como se mostró en el Capítulo~\ref{chap:diagnostico},
    el consumo real en los BS del Caribe supera esos valores entre un
    61\% y un 98\%.

    \item \textbf{Capacidad disponible en la red del OR:} Verificar que
    el circuito y el transformador más cercano al BS puedan absorber
    tanto la carga normalizada como la inyección del AGPE sin superar
    el límite del 50\% de la CREG 174 de 2021
    \parencite{ResolucionCREG174_2021}.

    \item \textbf{Revisión de convocatorias anteriores del PRONE:}
    Verificar las condiciones establecidas en convocatorias previas para
    identificar requisitos recurrentes, criterios de priorización y
    causales de descalificación que puedan anticiparse en la formulación.
\end{enumerate}

\textbf{Productos de la fase:}
\begin{itemize}
    \item Informe de diagnóstico integral del BS.
    \item Certificación de BS vigente (o constancia de gestión ante la
    alcaldía).
    \item Censo de usuarios con levantamiento técnico de la red existente.
    \item Datos de recurso solar del sitio.
    \item Informe de capacidad disponible en la red del OR.
\end{itemize}

\textbf{Criterio de avance:} Se puede pasar a la Fase 2 solo cuando
existen la certificación de BS vigente, el censo completo y los datos
de recurso solar. Sin la certificación, el proceso se detiene.

\textbf{Cuello de botella:} La certificación del BS depende de la
voluntad y los tiempos del ente territorial. Mitigación: iniciar las
gestiones seis meses antes de la convocatoria y mantener seguimiento
activo con la entidad territorial. Sin esta certificación, el proyecto
no puede avanzar a ninguna fase posterior.

\subsection{Fase 2 -- Diseño y dimensionamiento del proyecto}
\label{subsec:fase2}

\textbf{Objetivo:} Desarrollar el diseño completo de la red de
normalización y del sistema AGPE, y verificar que técnicamente son
compatibles entre sí y con la red del OR.

\textbf{Responsables:} OR (dirección técnica), contratista o consultora
de diseño, OR y CREG (verificación de capacidad de conexión).

\textbf{Actividades:}

\begin{enumerate}[label=\textbf{\arabic*.}, leftmargin=1.8cm, itemsep=6pt]

    \item \textbf{Diseño de la red eléctrica normalizada} conforme al
    RETIE, incluyendo los ítems de la Sección~\ref{subsec:disenio_electrico}.
    Se recomienda conductor concéntrico antifraude en BT, consistente
    con la experiencia documentada de ESSA en Santander
    \parencite{MendezFuentes2013}. El diseño debe cerrar con diagramas
    unifilares, planos de construcción y presupuesto de obra detallado.

    \item \textbf{Elaboración de memorias de cálculo y planos:} Diseños
    de ingeniería completos, memorias de cálculo y planos eléctricos y
    civiles, los cuales deben ser aportados a título gratuito al MME
    como parte del expediente.

    \item \textbf{Verificar la capacidad de conexión} consultando el
    micrositio de información de usuarios autogeneradores del OR (exigido
    por la CREG 174 de 2021). Confirmar que la potencia del AGPE no
    supera el 50\% de la capacidad disponible en el transformador o
    circuito \parencite{ResolucionCREG174_2021}.

    \item \textbf{Dimensionamiento del sistema AGPE} con base en el
    recurso solar del sitio, la demanda del BS normalizado y el límite
    del 50\%. Incluye selección de paneles, configuración de inversores
    (string o central) y sistema de medición bidireccional. Si el
    proyecto contempla almacenamiento, dimensionar el banco de baterías
    con la autonomía y profundidad de descarga requeridas. Debe incluirse
    el análisis de alternativas del sistema AGPE y los diagramas
    unifilares existentes y proyectados, así como el sistema de
    infraestructura de medición avanzada (AMI).

    \item \textbf{Definir quién es dueño del AGPE:} ¿El OR o la
    Comunidad Energética? Esta decisión determina la estructura legal
    del proyecto y el plan de AOM. Si se opta por la Comunidad
    Energética, iniciar su constitución conforme al Decreto 2236 de
    2023 \parencite{Decreto2236_2023}.

    \item \textbf{Plan de sostenibilidad:} Documento donde se determinen
    las estrategias para garantizar la sostenibilidad del proyecto a
    nivel social, ambiental y económico, incluyendo el plan de AOM para
    la red normalizada y el sistema AGPE.

    \item \textbf{Cronograma del proyecto:} Cronograma proyectado para
    la ejecución del proyecto y las actividades requeridas, alineado con
    los plazos de la convocatoria del PRONE.

    \item \textbf{Presupuesto integrado:} Un solo documento que
    consolide red de normalización y AGPE, con cantidades de obra,
    precios de mercado verificables y costo por usuario para los dos
    escenarios (solo red vs.\ red con AGPE). Incluye presupuesto a nivel
    global y unitario.
\end{enumerate}

\textbf{Productos de la fase:}
\begin{itemize}
    \item Diseño eléctrico completo de la red con memorias de cálculo y
    planos.
    \item Especificaciones y dimensionamiento del AGPE con análisis de
    alternativas.
    \item Verificación de capacidad de conexión.
    \item Presupuesto integrado con costo por usuario.
    \item Plan de sostenibilidad.
    \item Cronograma del proyecto.
    \item Decisión documentada sobre el modelo de propiedad de los
    activos.
\end{itemize}

\textbf{Criterio de avance:} El AGPE cumple el límite del 50\% y el
presupuesto es coherente con los rangos históricos del PRONE.

\textbf{Riesgos:} El más frecuente es que la capacidad disponible del
transformador o circuito existente no alcance para absorber la inyección
del AGPE, lo que obliga a incluir el refuerzo del punto de conexión en
el alcance — con el consiguiente aumento del presupuesto. Otro riesgo
real es que los requisitos de la convocatoria cambien entre el inicio
del diseño y la fecha de presentación: siempre revisar la resolución
vigente antes de cerrar el expediente.

\subsection{Fase 3 -- Estructuración y gestión del financiamiento}
\label{subsec:fase3}

\textbf{Objetivo:} Consolidar un expediente completo, coherente y
elegible, y presentarlo al MME dentro del plazo de la convocatoria para
su evaluación y aprobación por el CAPRONE.

\textbf{Responsables:} OR (consolidación y presentación), alcaldía
(certificación y soportes), consultora técnica (revisión de elegibilidad).

\textbf{Actividades:}

\begin{enumerate}[label=\textbf{\arabic*.}, leftmargin=1.8cm, itemsep=6pt]

    \item \textbf{Revisar la resolución de convocatoria vigente} antes
    de cerrar el expediente. Los requisitos cambian entre vigencias. La
    Resolución 40243 de 2025 introdujo cambios significativos al
    incorporar el AGPE como componente elegible
    \parencite{ResolucionMME40243_2025}.

    \item \textbf{Formular el proyecto en la MGA:} Presentar el proyecto
    en el formulario de la Metodología General Ajustada, herramienta
    diseñada por el DNP para la formulación y evaluación de proyectos
    de inversión pública.

    \item \textbf{Consolidar el expediente completo:} Integrar todos los
    documentos requeridos: certificación del BS, diagnóstico integral,
    diseño técnico completo (planos, memorias, unifilares),
    especificaciones del AGPE, verificación del límite del 50\%,
    presupuesto detallado, plan de gestión social, acuerdos comunitarios,
    análisis de impacto tarifario, plan de inversiones quinquenal,
    garantía de seriedad de la oferta, carta del ente territorial sobre
    alumbrado público, listado de usuarios beneficiados, certificaciones
    del OR y, si aplica, acto de constitución de la Comunidad Energética
    y documento de cofinanciación.

    \item \textbf{Suscribir el Plan de Acción de Responsabilidad Social}
    del MME, incluyendo las obligaciones asociadas a Comunidades
    Energéticas cuando el proyecto incorpore AGPE.

    \item \textbf{Revisión cruzada de elegibilidad:} Verificar ítem por
    ítem el cumplimiento de los requisitos de la convocatoria. Esta
    revisión es lo que en 2024 habría evitado que todos los proyectos
    presentados resultaran descalificados \parencite{MME2025}. Se
    recomienda que la haga un profesional externo al equipo formulador.

    \item \textbf{Radicar el expediente} ante el MME dentro del plazo,
    conservando la constancia del radicado.

    \item \textbf{Evaluación del CAPRONE:} Una vez radicado, el CAPRONE
    evalúa el proyecto conforme a los criterios de la convocatoria.
    Atender las solicitudes de información adicional del CAPRONE durante
    la evaluación, dentro de los plazos establecidos.
\end{enumerate}

\textbf{Productos de la fase:}
\begin{itemize}
    \item Expediente completo radicado ante el MME.
    \item Proyecto formulado en la MGA.
    \item Constancia de radicado.
    \item Plan de Acción de Responsabilidad Social suscrito.
    \item Respuestas documentadas al CAPRONE (si aplica).
\end{itemize}

\textbf{Riesgo principal:} Expedientes incompletos o con inconsistencias
entre los componentes técnico, financiero y social. Un profesional
externo con experiencia en evaluación de proyectos del PRONE que revise
el expediente antes de radicarlo puede evitar una descalificación
innecesaria.

\subsection{Fase 4 -- Ejecución, energización y sostenibilidad}
\label{subsec:fase4}

\textbf{Objetivo:} Ejecutar las obras, poner en servicio la
infraestructura y garantizar que los resultados se mantengan en el
largo plazo.

\textbf{Responsables:} OR (supervisión y AOM de la red), contratista
(obras), Comunidad Energética o OR (AOM del AGPE), alcaldía y entidades
de apoyo (seguimiento institucional).

\textbf{Actividades:}

\begin{enumerate}[label=\textbf{\arabic*.}, leftmargin=1.8cm, itemsep=6pt]

    \item \textbf{Contratar y ejecutar las obras} conforme al diseño
    aprobado. La interventoría es obligatoria y su costo no debe superar
    el 7\% de los costos directos del proyecto. Debe garantizarse la
    contratación de al menos el 50\% de personal local conforme a las
    certificaciones del OR.

    \item \textbf{Mantener activa la gestión social durante las obras.}
    El caso de Brisas del Venado documenta que la resistencia puede
    aparecer en cualquier momento de la ejecución, incluso cuando la
    socialización previa fue exitosa.

    \item \textbf{Pruebas de calidad y energización:} Verificar
    instalaciones conforme al RETIE, energizar la red y confirmar el
    correcto funcionamiento de los medidores individuales y del sistema
    de medición bidireccional del AGPE.

    \item \textbf{Puesta en operación del AGPE:} Verificar el
    funcionamiento del sistema fotovoltaico, la inyección o el
    autoconsumo (según la configuración) y la contabilización correcta
    en el medidor bidireccional. Si los activos se ceden a la Comunidad
    Energética, formalizar la entrega y capacitar a los gestores de la
    comunidad.

    \item \textbf{Apoyo a redes internas y punto de conexión AGPE:}
    Conforme a las obligaciones adquiridas por el OR, garantizar el
    apoyo a la adecuación de redes internas de los usuarios y la
    correcta habilitación del punto de conexión del sistema AGPE.

    \item \textbf{AOM de largo plazo:} El OR responde por la red de
    normalización de manera indefinida. Para el AGPE, la responsabilidad
    es del OR o de la Comunidad Energética según lo definido en la
    Fase 2. El plan de AOM debe prever inspecciones semestrales en el
    primer año, protocolo de respuesta ante fallas del inversor o los
    paneles (con directorio de proveedores con servicio técnico en la
    región), plan de limpieza de paneles — en el Caribe la salinidad
    y el polvo pueden reducir significativamente el rendimiento si no
    se limpia con regularidad — y seguimiento del recaudo con acciones
    correctivas cuando caiga por debajo de lo proyectado.

    \item \textbf{Evaluación de resultados a los 12 y 36 meses:}
    Porcentaje de recaudo, reducción de pérdidas, energía generada por
    el AGPE, reducción promedio de la factura del usuario, satisfacción
    comunitaria y número de usuarios que reincidieron en conexiones
    ilegales. Estos indicadores deben reportarse al MME conforme a las
    obligaciones de seguimiento de la resolución de adjudicación.
\end{enumerate}

\textbf{Productos de la fase:}
\begin{itemize}
    \item Acta de energización de la red y del AGPE.
    \item Acta de entrega de activos a la Comunidad Energética (si aplica).
    \item Informes de seguimiento a los 12 y 36 meses.
    \item Plan de AOM documentado y en ejecución.
\end{itemize}

\textbf{Riesgos principales:} La reincidencia en conexiones ilegales —
más frecuente en el primer año, cuando los usuarios aún están ajustando
sus hábitos de pago — y el deterioro prematuro del sistema AGPE por
falta de mantenimiento. El primero se mitiga con presencia periódica del
OR en el barrio. El segundo con un plan de AOM que tenga recursos reales
asignados, no solo intenciones.

\subsection{Fase Transversal -- Gestión social y comunitaria}
\label{subsec:fase_transversal}

\textbf{Objetivo:} Construir las condiciones sociales que hagan viable
la aceptación y la sostenibilidad del proyecto.

\textbf{Responsables:} OR (coordinación), trabajo social (implementación),
JAC y líderes comunitarios (interlocución), alcaldía y entidades de
apoyo (articulación institucional).

Esta fase no tiene un inicio y un fin discretos. Es transversal a las
cuatro fases secuenciales: empieza en el diagnóstico y no termina hasta
al menos un año después de la energización. Sus momentos más intensivos
son tres:

\textbf{Antes de las obras:}
\begin{itemize}
    \item Sesiones de socialización en lenguaje accesible: qué es la
    normalización, qué va a pasar en el barrio, cuánto va a costar el
    servicio, cómo funciona el AGPE y qué impacto tiene en la factura.
    La información clara sobre la tarifa es el factor que más reduce la
    resistencia inicial, según los 20 expertos consultados.
    \item Identificar y vincular activamente a los líderes con
    influencia real. Un líder que entiende y apoya el proyecto vale
    más que cualquier jornada de socialización masiva.
    \item Abordar explícitamente la desconfianza histórica hacia el OR.
    En algunos barrios ese trabajo toma más tiempo que la propia obra.
    \item Definir concertadamente el tratamiento de las deudas históricas.
    Sin condonación o refinanciación clara, la normalización empieza con
    una carga que muchas comunidades no van a aceptar
    \parencite{SerranoRueda2019}.
    \item Si hay AGPE y se prevé Comunidad Energética, iniciar la
    formación de la comunidad sobre sus derechos y responsabilidades
    como prestadora de servicio público.
\end{itemize}

\textbf{Durante las obras:}
\begin{itemize}
    \item Gestionar los conflictos que surjan: acceso a predios, cortes
    temporales, ruido, daños a bienes privados. Que la obra avance sin
    sorpresas para la comunidad.
    \item Capacitar a los usuarios en el uso seguro del nuevo servicio,
    lectura de la factura y acceso al subsidio FOES.
    \item Si hay Comunidad Energética, formalizar su constitución durante
    este periodo con asesoría jurídica y técnica.
\end{itemize}

\textbf{Post-normalización:}
\begin{itemize}
    \item Monitorear el recaudo durante los primeros seis meses con
    visitas a los sectores de menor pago.
    \item Ejecutar el plan de eficiencia energética: capacitación en
    uso racional y, si el presupuesto lo permite, el ``Cambiatón'' de
    luminarias que tan bien funcionó en Popayán
    \parencite{SerranoRueda2019}.
    \item Para proyectos con AGPE: capacitar en el monitoreo básico del
    sistema (limpieza de paneles, revisión visual de inversores, canales
    de reporte de fallas).
    \item Evaluaciones de satisfacción a los tres y doce meses.
\end{itemize}

\textbf{Productos de la fase transversal:}
\begin{itemize}
    \item Plan de gestión social documentado.
    \item Actas de sesiones de socialización.
    \item Acuerdo con el suscriptor comunitario o JAC.
    \item Decisión documentada sobre las deudas históricas.
    \item Acto de constitución de la Comunidad Energética (si aplica).
    \item Informes de seguimiento post-normalización.
\end{itemize}

\textbf{Riesgo principal:} La presencia de suscriptores comunitarios o
líderes con interés económico en mantener el esquema informal. Cobran
por las conexiones ilegales o administran el macromedidor y no quieren
perder esa posición. Identificarlos temprano y gestionar sus intereses
de manera explícita es clave para evitar que organicen la oposición. En
comunidades étnicas, los mecanismos de consulta y toma de decisiones
son diferentes y requieren tiempos distintos \parencite{DANE2018}.


% ─────────────────────────────────────────────────────────────
\section{Análisis de Replicabilidad y Recomendaciones}
\label{sec:replicabilidad}
% ─────────────────────────────────────────────────────────────

La metodología propuesta no fue diseñada para un solo barrio ni para
un único OR. Su valor radica en la posibilidad de ser aplicada de manera
sistemática en distintos contextos territoriales, siempre que se cumplan
las condiciones habilitantes del marco normativo y se adapten las
variables técnicas y sociales a las particularidades de cada caso. Esta
sección analiza los criterios de replicabilidad, identifica los vacíos
regulatorios que dificultan la implementación y formula recomendaciones
para superarlos.

\subsection{Criterios de replicabilidad}
\label{subsec:criterios_replicabilidad}

\textbf{Normativa transversal:} La metodología se fundamenta en
normativa de alcance nacional — el DUR 1073 de 2015, la CREG 174 de
2021, el Decreto 2236 de 2023, la Resolución 40243 de 2025, entre
otras —, lo que permite su aplicación en cualquier departamento del país
donde existan BS registrados en el SUI. No depende de normativa local
específica, aunque sí requiere la voluntad del ente territorial para
expedir la certificación del BS.

\textbf{Escalabilidad financiera:} El esquema de financiamiento a
través del PRONE es el mismo para todos los OR del SIN. Los requisitos
de elegibilidad, los criterios de priorización y los mecanismos de
asignación de recursos están definidos en las resoluciones de
convocatoria del MME y aplican de manera uniforme. Esto permite que
cualquier OR con capacidad de formulación —o que contrate esa capacidad—
pueda acceder al programa.

\textbf{Adaptabilidad técnica:} Las variables de diseño del sistema
AGPE (recurso solar, demanda, capacidad de la red) son específicas de
cada sitio, pero la secuencia metodológica para determinarlas es la
misma. Un OR en el Caribe y un OR en el Pacífico siguen la misma hoja
de ruta; lo que cambia es el resultado del dimensionamiento, no el
proceso para llegar a él.

\textbf{Componente social transferible:} Las estrategias de gestión
social —socialización, vinculación de líderes, tratamiento de deudas,
formación de Comunidades Energéticas— son aplicables en cualquier BS,
con las adaptaciones culturales necesarias. La experiencia de Popayán
\parencite{SerranoRueda2019} y los resultados de la consulta a expertos
muestran que los factores de éxito social son consistentes
independientemente de la región.

\subsection{Vacíos regulatorios identificados}
\label{subsec:vacios_regulatorios}

El análisis normativo del Capítulo~\ref{chap:normativo} y la aplicación
práctica de los requisitos del PRONE al diseño de la metodología
revelaron cuatro vacíos que hoy dificultan la implementación efectiva
del AGPE dentro del esquema del programa.

\textbf{Vacío 1 -- Sin incentivos para que la Comunidad Energética
gestione el AGPE.} El Decreto 2236 de 2023 habilita la cesión de
activos a las Comunidades Energéticas, pero no establece ningún
mecanismo de remuneración por esa responsabilidad. Pedirle a una
comunidad con capacidades técnicas y administrativas limitadas que
gestione un sistema fotovoltaico sin ningún incentivo económico es
poco realista \parencite{Decreto2236_2023}.

\textbf{Vacío 2 -- El subsidio FOES no está articulado con el AGPE.}
No existe ninguna disposición normativa que defina cómo se calcula el
FOES para usuarios normalizados que cuentan con generación solar. El
usuario reduce su consumo neto de la red, pero el subsidio se calcula
sobre el consumo total. Esa incoherencia fue señalada por el experto
E11 en la consulta: puede hacer que el impacto en la factura sea más
crítico de lo esperado.

\textbf{Vacío 3 -- Sin criterios de priorización específicos para
proyectos con AGPE.} Las convocatorias usan criterios generales para
todos los proyectos, sin reconocer que los proyectos con AGPE tienen
mayor complejidad de formulación y mayor costo unitario que los
convencionales \parencite{ResolucionMME40243_2025}.

\textbf{Vacío 4 -- El PRONE no exige demostrar sostenibilidad
post-normalización.} El programa financia la ejecución, pero no tiene
mecanismos que obliguen a los OR a demostrar que los resultados se
mantienen en el tiempo. Sin ese componente, se corre el riesgo de
financiar proyectos que funcionan bien en la foto inicial y se
deterioran en los años siguientes.

\subsection{Recomendaciones}
\label{subsec:recomendaciones}

A partir de los vacíos identificados y del análisis de replicabilidad,
se formulan las siguientes recomendaciones dirigidas a los actores del
proceso:

\begin{enumerate}[label=\textbf{\arabic*.}, leftmargin=1.8cm, itemsep=6pt]

    \item \textbf{Definir un esquema de remuneración para las Comunidades
    Energéticas.} El MME y la CREG deberían establecer un mecanismo de
    remuneración estructurado como un porcentaje de los excedentes
    entregados a la red o como reconocimiento tarifario dentro del FOES,
    de manera que la gestión del AGPE sea económicamente viable para la
    comunidad.

    \item \textbf{Articular el subsidio FOES con el AGPE.} El MME, la
    CREG y la UPME deben desarrollar una disposición que defina el
    cálculo del FOES para usuarios con AGPE, proporcional al consumo
    neto de la red, evitando que el beneficio de la generación solar se
    diluya por una formulación tarifaria incompatible.

    \item \textbf{Incluir criterios diferenciados de priorización para
    proyectos con AGPE.} En futuras resoluciones de convocatoria, se
    recomienda valorar el recurso solar disponible, el impacto en la
    factura del usuario, el potencial de reducción de pérdidas y la
    solidez del plan de AOM como criterios específicos de priorización.

    \item \textbf{Incorporar informes de impacto obligatorios
    post-normalización.} Se recomienda que el PRONE exija informes de
    seguimiento a los 24 y 48 meses de la energización, y que vincule
    el desempeño en proyectos anteriores a la elegibilidad para futuras
    convocatorias.

    \item \textbf{Fortalecer la capacidad de formulación de los OR.}
    Muchos OR con BS en su área de influencia no cuentan con la
    capacidad interna para formular proyectos que cumplan todos los
    requisitos del PRONE. Se recomienda que el MME, en articulación con
    la UPME, desarrolle programas de asistencia técnica y guías de
    formulación que reduzcan las barreras de entrada al programa.

    \item \textbf{Estandarizar los requisitos entre convocatorias.}
    La variabilidad de requisitos entre convocatorias dificulta la
    preparación anticipada de los expedientes. Establecer un núcleo
    estable de requisitos con variaciones marginales entre vigencias
    permitiría a los OR iniciar la formulación con mayor certeza.
\end{enumerate}